% GENERATED FILE. Edit canonical lessons, then run npm run generate:books.
% canonical-source-hash: fnv1a64:5dacd9238e1df815
% canonical-lessons: ML-C523-manasikavastha, ML-C523-paka, ML-C523-pirupirukkuka, ML-C523-kusumpu, ML-C523-aradhikkuka

\chapter{Grudges and Memories}
\label{ch:ml-grudges-and-memories}
\hlchaptermodality{523}

\begin{chapteropening}
\textbf{By the end of this chapter:} \emph{I can say a state of mind, a grudge, to mutter, petty spite and to admire.}

Complete the last lesson of chapter 523: I can say a state of mind, a grudge, to mutter, petty spite and to admire.
\end{chapteropening}

\section[\texorpdfstring{\ml{മാനസികാവസ്ഥ} (mānasikāvastha)}{മാനസികാവസ്ഥ (mānasikāvastha)}]{\ml{മാനസികാവസ്ഥ} (mānasikāvastha) — a state of mind; a mood}

\label{lesson:ML-C523-manasikavastha}



\begin{quote}
Before the new one: say the Malayalam for freshness, then the Malayalam for to get fed up.
\end{quote}

\subsection*{You'll want to know: \ml{മാനസികാവസ്ഥ}}
\textbf{\ml{മാനസികാവസ്ഥ}} — \emph{mānasikāvastha} — \textquotedblleft{}a state of mind; a mood\textquotedblright{}.

\begin{quote}
\ml{പരീക്ഷയ്ക്ക്} \ml{മുൻപ്} \ml{അവന്റെ} \ml{മാനസികാവസ്ഥ} \ml{മോശമായിരുന്നു}.
\end{quote}

\begin{quote}
\emph{parīkṣaykkŭ munpŭ avanṟe mānasikāvastha mōśamāyirunnu.}
\end{quote}

\begin{quote}
Before the exam he was in a bad state of mind.
\end{quote}

\ml{മാനസിക} of the mind + \ml{അവസ്ഥ} a state: the two vowels join into one long \ml{ആ}.

\begin{cousinweb}
\textbf{\ml{മാനസികാവസ്ഥ}} (\emph{mānasikāvastha}) is a learned compound built from Sanskrit \emph{mānasika}, \textquotedblleft{}of the mind\textquotedblright{} (Monier-Williams), and \emph{avasthā}, \textquotedblleft{}a state, a condition\textquotedblright{} (CDIAL 872).
\end{cousinweb}

\begin{grammarlens}[title={What you've built}]
One more word for describing what you lived and felt.
\end{grammarlens}

\subsection*{Your turn}
\begin{itemize}
\raggedright
  \item \emph{Say it:} \emph{mānasikāvastha}
  \item \emph{Say it:} \emph{mānasikāvastha}, once more
  \item \emph{Say it:} \emph{parīkṣaykkŭ munpŭ avanṟe mānasikāvastha mōśamāyirunnu.} — Before the exam he was in a bad state of mind.
  \item \emph{Recall:} say the Malayalam for freshness, then the Malayalam for to get fed up, then say \emph{mānasikāvastha} again
\end{itemize}

\begin{tcolorbox}[breakable,colback=teal!4,colframe=teal!35!black,title={Before you move on}]
What does \ml{മാനസികാവസ്ഥ} mean? (\textquotedblleft{}A state of mind; a mood\textquotedblright{}.) Say it once more.
\end{tcolorbox}
\section[\texorpdfstring{\ml{പക} (paka)}{പക (paka)}]{\ml{പക} (paka) — a grudge; vengefulness}

\label{lesson:ML-C523-paka}



\begin{quote}
Before the new one: say the Malayalam for to get fed up, then the Malayalam for a state of mind.
\end{quote}

\subsection*{You'll want to know: \ml{പക}}
\textbf{\ml{പക}} — \emph{paka} — \textquotedblleft{}a grudge; vengefulness\textquotedblright{}.

\begin{quote}
\ml{പഴയ} \ml{വഴക്കിന്റെ} \ml{പക} \ml{അവൻ} \ml{മറന്നില്ല}.
\end{quote}

\begin{quote}
\emph{paḻaya vaḻakkinṟe paka avan maṟannilla.}
\end{quote}

\begin{quote}
He did not forget the grudge from the old quarrel.
\end{quote}

\ml{പക} \ml{വീട്ടുക} is to take revenge, literally to pay off a grudge.

\begin{cousinweb}
\textbf{\ml{പക}} (\emph{paka}) is native Dravidian: DEDR 3808 lists it beside verbs for dividing and sharing, with the sense of separation and enmity (DEDR 3808).
\end{cousinweb}

\begin{grammarlens}[title={What you've built}]
One more word for describing what you lived and felt.
\end{grammarlens}

\subsection*{Your turn}
\begin{itemize}
\raggedright
  \item \emph{Say it:} \emph{paka}
  \item \emph{Say it:} \emph{paka}, once more
  \item \emph{Say it:} \emph{paḻaya vaḻakkinṟe paka avan maṟannilla.} — He did not forget the grudge from the old quarrel.
  \item \emph{Recall:} say the Malayalam for to get fed up, then the Malayalam for a state of mind, then say \emph{paka} again
\end{itemize}

\begin{tcolorbox}[breakable,colback=teal!4,colframe=teal!35!black,title={Before you move on}]
What does \ml{പക} mean? (\textquotedblleft{}A grudge; vengefulness\textquotedblright{}.) Say it once more.
\end{tcolorbox}
\section[\texorpdfstring{\ml{പിറുപിറുക്കുക} (piṟupiṟukkuka)}{പിറുപിറുക്കുക (piṟupiṟukkuka)}]{\ml{പിറുപിറുക്കുക} (piṟupiṟukkuka) — to mutter}

\label{lesson:ML-C523-pirupirukkuka}



\begin{quote}
Before the new one: say the Malayalam for a state of mind, then the Malayalam for a grudge.
\end{quote}

\subsection*{You'll want to know: \ml{പിറുപിറുക്കുക}}
\textbf{\ml{പിറുപിറുക്കുക}} — \emph{piṟupiṟukkuka} — \textquotedblleft{}to mutter\textquotedblright{}.

\begin{quote}
\ml{മുത്തശ്ശി} \ml{എന്തോ} \ml{പിറുപിറുത്തുകൊണ്ട്} \ml{അടുക്കളയിലേക്ക്} \ml{പോയി}.
\end{quote}

\begin{quote}
\emph{muttaśśi entō piṟupiṟuttukoṇṭŭ aṭukkaḷayilēkkŭ pōyi.}
\end{quote}

\begin{quote}
Muttering something, grandmother went into the kitchen.
\end{quote}

Past: \ml{പിറുപിറുത്തു}, \emph{piṟupiṟuttu}.

Talking to oneself in a low voice; the past is \ml{പിറുപിറുത്തു}.

\begin{cousinweb}
\textbf{\ml{പിറുപിറുക്കുക}} (\emph{piṟupiṟukkuka}) is a doubled sound-word like \ml{മുറുമുറുക്കുക}; DEDR has no entry that clearly holds it, so its origin is not settled.
\end{cousinweb}

\begin{grammarlens}[title={What you've built}]
One more word for describing what you lived and felt.
\end{grammarlens}

\subsection*{Your turn}
\begin{itemize}
\raggedright
  \item \emph{Say it:} \emph{piṟupiṟukkuka}
  \item \emph{Say it:} \emph{piṟupiṟukkuka}, once more
  \item \emph{Say it:} \emph{muttaśśi entō piṟupiṟuttukoṇṭŭ aṭukkaḷayilēkkŭ pōyi.} — Muttering something, grandmother went into the kitchen.
  \item \emph{Recall:} say the Malayalam for a state of mind, then the Malayalam for a grudge, then say \emph{piṟupiṟukkuka} again
\end{itemize}

\begin{tcolorbox}[breakable,colback=teal!4,colframe=teal!35!black,title={Before you move on}]
What does \ml{പിറുപിറുക്കുക} mean? (\textquotedblleft{}To mutter\textquotedblright{}.) Say it once more.
\end{tcolorbox}
\section[\texorpdfstring{\ml{കുശുമ്പ്} (kuśumpŭ)}{കുശുമ്പ് (kuśumpŭ)}]{\ml{കുശുമ്പ്} (kuśumpŭ) — petty spite; envy}

\label{lesson:ML-C523-kusumpu}



\begin{quote}
Before the new one: say the Malayalam for a grudge, then the Malayalam for to mutter.
\end{quote}

\subsection*{You'll want to know: \ml{കുശുമ്പ്}}
\textbf{\ml{കുശുമ്പ്}} — \emph{kuśumpŭ} — \textquotedblleft{}petty spite; envy\textquotedblright{}.

\begin{quote}
\ml{അനിയന്റെ} \ml{പുതിയ} \ml{സൈക്കിൾ} \ml{കണ്ട്} \ml{അവന്} \ml{കുശുമ്പ്} \ml{തോന്നി}.
\end{quote}

\begin{quote}
\emph{aniyanṟe putiya saikkiḷ kaṇṭŭ avanŭ kuśumpŭ tōnni.}
\end{quote}

\begin{quote}
He felt a pang of envy at his little brother's new bicycle.
\end{quote}

Spoken and a little playful; \ml{അസൂയ} is the general word for envy.

\begin{cousinweb}
\textbf{\ml{കുശുമ്പ്}} (\emph{kuśumpŭ}) is a spoken word; DEDR has no clear entry for it, so its origin is not settled.
\end{cousinweb}

\begin{grammarlens}[title={What you've built}]
One more word for describing what you lived and felt.
\end{grammarlens}

\subsection*{Your turn}
\begin{itemize}
\raggedright
  \item \emph{Say it:} \emph{kuśumpŭ}
  \item \emph{Say it:} \emph{kuśumpŭ}, once more
  \item \emph{Say it:} \emph{aniyanṟe putiya saikkiḷ kaṇṭŭ avanŭ kuśumpŭ tōnni.} — He felt a pang of envy at his little brother's new bicycle.
  \item \emph{Recall:} say the Malayalam for a grudge, then the Malayalam for to mutter, then say \emph{kuśumpŭ} again
\end{itemize}

\begin{tcolorbox}[breakable,colback=teal!4,colframe=teal!35!black,title={Before you move on}]
What does \ml{കുശുമ്പ്} mean? (\textquotedblleft{}Petty spite; envy\textquotedblright{}.) Say it once more.
\end{tcolorbox}
\section[\texorpdfstring{\ml{ആരാധിക്കുക} (ārādhikkuka)}{ആരാധിക്കുക (ārādhikkuka)}]{\ml{ആരാധിക്കുക} (ārādhikkuka) — to admire; to adore}

\label{lesson:ML-C523-aradhikkuka}



\begin{quote}
Before the new one: say the Malayalam for to mutter, then the Malayalam for petty spite.
\end{quote}

\subsection*{You'll want to know: \ml{ആരാധിക്കുക}}
\textbf{\ml{ആരാധിക്കുക}} — \emph{ārādhikkuka} — \textquotedblleft{}to admire; to adore\textquotedblright{}.

\begin{quote}
\ml{എന്റെ} \ml{അനിയത്തി} \ml{ആ} \ml{ഗായകനെ} \ml{ആരാധിക്കുന്നു}.
\end{quote}

\begin{quote}
\emph{enṟe aniyatti ā gāyakane ārādhikkunnu.}
\end{quote}

\begin{quote}
My little sister adores that singer.
\end{quote}

Past: \ml{ആരാധിച്ചു}, \emph{ārādhiccu}.

The one admired takes -\ml{എ}. It also means to worship; a fan is \ml{ആരാധകൻ}.

\begin{cousinweb}
\textbf{\ml{ആരാധിക്കുക}} (\emph{ārādhikkuka}) is built from Sanskrit \emph{ārādhana}, the noun of \emph{ārādhayati}, \textquotedblleft{}honours, wins over\textquotedblright{} (CDIAL 1316), with the Malayalam verb ending -\ml{ിക്കുക}.
\end{cousinweb}

\begin{grammarlens}[title={What you've built}]
One more word for describing what you lived and felt.
\end{grammarlens}

\subsection*{Your turn}
\begin{itemize}
\raggedright
  \item \emph{Say it:} \emph{ārādhikkuka}
  \item \emph{Say it:} \emph{ārādhikkuka}, once more
  \item \emph{Say it:} \emph{enṟe aniyatti ā gāyakane ārādhikkunnu.} — My little sister adores that singer.
  \item \emph{Recall:} say the Malayalam for to mutter, then the Malayalam for petty spite, then say \emph{ārādhikkuka} again
\end{itemize}

\begin{tcolorbox}[breakable,colback=teal!4,colframe=teal!35!black,title={Before you move on}]
What does \ml{ആരാധിക്കുക} mean? (\textquotedblleft{}To admire; to adore\textquotedblright{}.) Say it once more.
\end{tcolorbox}
